\section{Related Work}
\label{sec:related}

\paragraph{Transaction chopping and its descendants.}
Chopping, the SC-graph criterion and the non-removability of cycles under
refinement are due to Shasha et al.~\cite{shasha1995chopping}; the
serializability theory underneath is
classical~\cite{bernstein1987,eswaran1976}. Later systems combine a static
analysis of pieces with runtime mechanisms to admit executions that a static
no-cycle test would reject: transaction chains order pieces across
sites~\cite{zhang2013chains}, Callas groups transactions and applies a
different mechanism within and across groups~\cite{xie2015callas}, and IC3
tracks dependencies between pieces at run time~\cite{wang2016ic3}. Chopping
has also been studied under parallel snapshot isolation, where the permitted
observers change what is safe~\cite{cerone2015psi}. All of these assume the
transaction programs are known when the analysis runs. Our criterion is
deliberately more conservative than the runtime-assisted ones and
correspondingly cheaper; what we add is the case they do not consider, a
program that changes after part of it has committed.

\paragraph{Deadlock handling.}
\ccs applies wound-wait~\cite{rosenkrantz1978} inside pieces, with the
launch-time order of the workflows as the timestamp. The scheme is old; the
observation that makes it useful here is that agents tolerate redo better than
repair, which is an empirical property of models rather than of locks.
Optimistic control~\cite{kung1981occ} and saga
compensation~\cite{garciamolina1987sagas} are the basis of our baselines.

\paragraph{Concurrency control for agents.}
CoAgent~\cite{lyu2026coagent} orders agents at launch and repairs conflicts
semantically; it is our main point of comparison, and its ordering is what
\ccs reuses for locks. Where CoAgent assumes a notified agent heals itself,
we measure how often it does and offer a design in which the corresponding
failure leaves state clean. Atomix~\cite{mohammadi2026atomix} provides
transactional tool use with declared effects and leaves the choice of
transaction granularity to the integration layer; certified chopping is a way
of making that choice. OptiMA~\cite{calikyilmaz2025optima} schedules
multi-agent transactions that it treats as indivisible.
SagaLLM~\cite{chang2025sagallm} applies compensation-based decomposition to
agent planning and explicitly relaxes isolation, and
Mnemosyne~\cite{chang2026mnemosyne} admits untrusted agent proposals through
a deterministic gate, a stance our fail-closed enforcement shares.
Concurrency anomalies in multi-agent systems have been characterized
formally~\cite{khan2026verified}. Automatic reconstruction of an agent's read
set~\cite{khan2026sbus} and adaptive control for agent-generated
transactions~\cite{zhou2026atcc} are complementary: the first could supply
the footprints we take as given.

\paragraph{Agent benchmarks.}
Our contended cells draw their operations from agent benchmarks for cloud
operations~\cite{chen2025aiopslab} and workplace
tasks~\cite{styles2024workbench}; $\tau$-bench~\cite{yao2024taubench} is a
further candidate. These benchmarks run tasks in isolation and therefore do
not, by themselves, create the contention this problem is about; adapting
them into concurrent versions is future work.

What we did not find in this body of work is a treatment of plan revision
over an already-committed prefix, with an admissibility characterization and
an impossibility result. That is a statement about what our search found, not
a priority claim.
